\section{Calibration design}\label{sec:calib}
Because the ambient PFOS record is a sparse set of per-station aggregate concentrations rather than continuous time series, calibration is posed as a \emph{spatial} matching problem: the model must reproduce the pattern of representative concentration across the $20$ gauged reaches ($29$ stations, station-maximum observed per reach), including the Rogue mainstem downstream gradient. For each station channel the modeled representative concentration is taken as the flow-weighted mean of daily outflow---total exported PFAS mass divided by total outflow volume---which is the statistic most directly comparable to grab samples collected across a range of flow conditions. Because PFAS concentrations span orders of magnitude across the network, the objective is evaluated in $\log_{10}$ space; we report the log-space Nash--Sutcliffe efficiency, root-mean-square error (in dex), and percent bias across all paired station--channel points, and minimize the log-space root-mean-square error.

The calibration is deliberately parsimonious, which both manages equifinality and exploits a structural property of the problem. In-stream concentration is, to first order, linear in the soil-loading magnitude: doubling the initial soil PFAS doubles the load delivered to runoff and therefore the reach concentration. The model therefore exposes a single global soil-loading scale factor (and, secondarily, a multiplier on the in-stream partition coefficient) applied uniformly across all HRUs, rather than independent per-reach parameters. This makes the magnitude calibration nearly analytic---one evaluation of the uncalibrated model fixes the scale factor to first order---and confines the calibrated degrees of freedom to physically grouped, class-based quantities. The concentrated Wolverine source is represented separately, as an optional per-reach point load in the lower river, so that the distributed-background and point-source contributions remain identifiable rather than absorbed into a single fitted field. Streamflow is calibrated independently against the Rockford gauge; we report the PFAS calibration on the uncalibrated-flow network because the spatial flow pattern that governs the concentration gradient is set by drainage area and is robust to flow calibration, which rescales concentration approximately uniformly and is re-absorbed by the soil-loading factor.

\subsection{Uncertainty analysis}
We quantify predictive uncertainty with a Monte-Carlo forward ensemble over the five parameters that govern PFAS fate, following the parameter-range propagation used for the prior distributed PFAS model \citep{rafiei2023}. A Latin-hypercube design of $120$ members samples uniform priors on the global soil-loading multiplier ($[0.05,0.22]$, bracketing the calibrated $0.11$), the in-stream organic-carbon partition multiplier ($[0.5,2.0]$), the Langmuir air--water-interface affinity $K_L$ ($[0.07,0.27]$) and maximum surface excess $\Gamma_{\max}$ ($[1500,3500]$), and the runoff/percolation partition ($[0.1,0.4]$); the full design is given in Supplementary Table~S2. Each member is a complete forward run of the engine on the Rogue, yielding a per-reach distribution of in-stream PFOS from which we form the $5$th--$95$th-percentile predictive envelope. Envelope quality is summarized by the $P$-factor (fraction of observations enclosed) and $R$-factor (mean envelope width normalized by the observed standard deviation), the standard SWAT calibration-uncertainty diagnostics. Parameter influence is ranked by the Spearman rank correlation between each sampled parameter and the reach-mean in-stream concentration. The $120$ forward runs were executed in parallel on a single $32$-vCPU cloud instance (AWS EC2 \texttt{c7i.8xlarge} spot); the full ensemble completed in $14.4$~min of compute ($15$~min $05$~s end to end including provisioning), at a mean of $172$~s per member.
